\section{Related Work}

This section organizes related work by both platform proximity and decision-layer proximity.
The paper does not address liquid sloshing in isolation, nor does it address generic mobile-robot local planning in isolation.
Instead, it focuses on their intersection at the online local-planning layer of a standard wheeled mobile base, where the planner must output executable commands at runtime while propagating the dynamic state of the transported liquid over the prediction horizon.

Three groups of prior work are most relevant.
The first group consists of ordinary mobile-robot local planners and trajectory optimization methods, which operate at a similar deployment layer but generally do not propagate liquid states.
The second group includes mobile-base liquid transport and closely related anti-slosh studies, which share the physical task but often operate through fixed paths, velocity profiles, offline trajectory optimization, tracking control, or special platform mechanisms.
The third group covers cross-platform slosh modeling, anti-slosh trajectory generation, and predictive control, which provide modeling and control foundations while clarifying the novelty boundary of this work.

\subsection{Online Local Planning and Trajectory Optimization for Mobile Robots}

Mobile-robot local planners address short-horizon motion generation at runtime and must output executable base commands under kinematic, dynamic, and environmental constraints.
Classical DWA, TEB, and \MpcLocalPlanner{} represent important families of velocity-space search, time-parameterized trajectory optimization, and receding-horizon nonlinear optimal control \cite{D15_Fox1997,D16_Rosmann2017,D17_Rosmann2021}.
\MPCC{} has also become a useful backbone for progress-maximizing path following and local planning \cite{D19_Liniger2015,D20_Brito2019}.
Regulated pure pursuit, dynamic-window variants, LT-DWA, MPPI, differential-drive trajectory optimization, and jerk-limited trajectory generation further extend local planning in terms of tracking, velocity and acceleration limits, dynamic feasibility, risk handling, and smoothness \cite{D01_Macenski2023,D22_Ohnishi2026,D02_Jian2023,D06_Williams2017,D07_Trevisan2025,D03_Zhang2025,D13_Berscheid2021}.

These methods are close to the system interface considered here because they serve runtime local decision making and ultimately produce executable base commands.
Local-planner benchmarks also typically evaluate safety, efficiency, smoothness, path tracking, and computation time \cite{D21_Wen2021}.
They therefore provide decision-layer context for \SMPCC{}, although the current fixed-path study does not rank global or obstacle-oriented planners.

However, these planners usually propagate robot pose, velocity, path error, obstacle distance, control variation, or risk variables; they do not include the modal displacement and velocity of the liquid inside the transported container.
For liquid transport, reducing command variation or improving geometric smoothness is not equivalent to suppressing sloshing.
Two command sequences may have similar velocity, acceleration, and jerk profiles but different future liquid responses because they couple differently with the current liquid phase.
Ordinary local planners clarify the same-layer novelty boundary, but they do not directly represent liquid dynamic memory.

\subsection{Mobile-Base Liquid Transport and Anti-Slosh Methods}

Mobile-base liquid transport is the closest physical task to this paper.
Hamaguchi and collaborators studied path design, velocity profiles, input shaping, and tracking control for wheeled mobile robots or carts transporting cylindrical liquid containers \cite{B04_Hamaguchi2002,B03_Hamaguchi2004,B02_Hamaguchi2005}.
These studies demonstrated that mobile-platform motion excites liquid sloshing and that motion planning and control benefit from explicit consideration of the liquid response.
Subsequent work also introduced active vibration reducers, parallel-link mechanisms, and omnidirectional mobile platforms \cite{B05_Hamaguchi2019,B06_Hamaguchi2018,B07_Nguyen2024,B08_Yuan2020}, showing that anti-slosh transport can be achieved through specialized hardware or platform designs as well as software planning.

Recent work has further incorporated liquid models into mobile-robot or mobile-container optimization and control.
Lim et al.\ approximate liquid sloshing by spherical-pendulum dynamics and solve a two-dimensional trajectory optimization problem to reduce sloshing \cite{B01_Lim2024}.
Nguyen Viet et al.\ combine a nonlinear equivalent mass-spring-damper model, time-optimal planning, tracking control, and output constraints for container transport under disturbances \cite{B11_Viet2025}.
Prabakaran et al.\ study slosh-aware trajectory control and \MPC{} tracking control on a reconfigurable stair-climbing service robot \cite{B12_Prabakaran2026}.

These studies establish the value of explicit liquid modeling and optimized motion generation.
Their main difference from this work is the decision layer.
The Hamaguchi line of work is closer to path design, velocity-profile design, input shaping, and tracking along a prescribed route.
Active suppression mechanisms rely on additional hardware or special platforms.
The approach of Lim et al.\ is closer to whole-trajectory offline optimization.
Nguyen Viet et al.\ use a planning-plus-tracking structure, and Prabakaran et al.\ focus on a special service-robot platform and trajectory tracking.
In contrast, \SMPCC{} inserts a low-order slosh state directly into an online \MPCC{} local-planning problem for a standard wheeled mobile base and resolves a short-horizon optimization problem at every control cycle.

\begin{table}[t]
  \centering
  \scriptsize
  \setlength{\tabcolsep}{3pt}
  \caption{Decision-layer comparison with mobile-base liquid-transport studies.}
  \label{tab:mobile_base_related_work_framework}
  \begin{tabular}{p{0.23\linewidth} p{0.29\linewidth} p{0.36\linewidth}}
    \toprule
    Type & Representative work & Main difference from this paper \\
    \midrule
    Fixed path, velocity profile, and input shaping & Hamaguchi series \cite{B04_Hamaguchi2002,B03_Hamaguchi2004,B02_Hamaguchi2005} & Suppresses sloshing mainly through predesigned paths or speed profiles rather than a path-progress \MPCC{} solved at every control cycle. \\
    Active suppression hardware & Mobile carts or omnidirectional platforms with active reducers \cite{B05_Hamaguchi2019,B06_Hamaguchi2018} & Relies on additional mechanisms or special platforms rather than a software local planner for a standard wheeled base. \\
    Offline liquid-constrained trajectory optimization & Lim et al. \cite{B01_Lim2024} & Uses explicit liquid dynamics but is closer to whole-trajectory optimization than runtime local planning. \\
    Time-optimal planning and tracking control & Nguyen Viet et al. \cite{B11_Viet2025} & Uses reference planning and robust tracking rather than directly producing local-planner base commands. \\
    Slosh-aware tracking on special platforms & Prabakaran et al. \cite{B12_Prabakaran2026} & Demonstrates slosh-aware \MPC{} or tracking control, but with different platform, task interface, and control layer. \\
    \bottomrule
  \end{tabular}
\end{table}

Thus, the mobile-base liquid-transport problem is not unexplored.
The gap addressed here is the insertion of liquid dynamic memory into the online local-planning loop of a standard mobile base.
\SMPCC{} defines this problem layer by selecting path progress, base control, and predicted liquid response in the same receding-horizon optimization problem.

\subsection{Cross-Platform Slosh Modeling, Trajectory Generation, and Predictive Control}

Liquid sloshing also appears in manipulators, SCARA systems, liquid-transfer machines, tank vehicles, ships, and specialized mechanisms.
High-fidelity VOF, FEM, CFD, and digital-twin models can describe complex free-surface behavior, fluid-structure coupling, and nonlinear effects \cite{A03_Schrgenhumer2019,A04_Elahi2015,A05_Nicolsen2017,A06_Moya2020}.
However, their state dimension and computational cost are generally unsuitable for direct embedding in a real-time mobile-robot local planner.
Equivalent-pendulum, mass-spring-damper, and low-order modal models instead provide a compact state representation of dominant slosh dynamics, making them more appropriate as prediction models for online optimization \cite{A02_Guagliumi2021,A01_Leva2021}.

Input shaping, predictive control, and anti-slosh trajectory generation have shown that explicit liquid models can improve motion generation and control \cite{A07_Singer1990,A12_Okatsuka2011,C05_Okatsuka2012,C01_Ferrari2026,C02_Muchacho2022,C03_Moriello2018,C04_Kuriyama2008,C06_Viet2024}.
Ferrari et al.\ propose offline time-optimal anti-slosh trajectory planning for multi-container transport on a 4D SCARA system and compare a low-order model height with video-extracted experimental liquid motion \cite{C01_Ferrari2026}.
That work provides a useful model-to-experiment validation perspective, but differs in platform, grasping configuration, optimization layer, and online interface.
Accordingly, this paper does not claim novelty as the first use of a slosh state in \MPC{}; it emphasizes the integration of a low-order slosh state into the online local-planning layer of a standard mobile base.

External sensing, machine vision, and liquid-level detection studies further show that real liquid motion can be evaluated using independent observations \cite{A08_Weaver2024,A09_Bobovnik2021,A10_Shao2026,A11_Ren2024}.
This motivates the distinction used throughout the paper between internal model proxies and measured liquid-surface indicators.
Cross-platform anti-slosh studies are therefore treated primarily as modeling and evaluation references unless their platform, input interface, and decision layer can be reproduced under the same experimental conditions.

\subsection{Positioning of \SMPCC{}}

In summary, prior work has established foundations in mobile-robot local planning, mobile-base liquid transport, low-order slosh modeling, and anti-slosh predictive control.
Ordinary local planners operate at a similar deployment layer but usually lack liquid-state prediction.
Mobile-base anti-slosh methods share the physical task but often operate through fixed paths, speed profiles, offline optimization, tracking control, or special platforms.
Cross-platform anti-slosh methods provide modeling and control foundations but do not generally constitute same-layer experimental baselines.

This paper is positioned at their intersection: online slosh-aware \MPCC{} local planning for a standard wheeled mobile base.
Rather than claiming the first slosh-aware \MPC{} method, \SMPCC{} integrates low-order slosh modal states, path progress, base-control constraints, control smoothness, and predicted liquid response into one receding-horizon local-planning problem. Its empirical claims remain limited to matched internal \MPCC{} variants, completion-time control, geometry-based container transfer, controlled liquid-phase planning, propagated-state replay, and shared execution-layer and runtime audits.
